\subsection{The product of two measures}

THEOREM 1. --- Let X and Y be two locally compact spaces, \(\lambda\) a measure on X, \(\mu\) a measure on Y; there exists one and only one measure \(\nu\) on X \(\times\) Y such that, for every function \(g \in \mathcal{K}(\mathrm{X};\mathbf{C})\) and every function \(h \in \mathcal{K}(\mathrm{Y};\mathbf{C})\) ,

\[
\int g (x) h (y) d \nu (x, y) = \left(\int g (x) d \lambda (x)\right) \left(\int h (y) d \mu (y)\right).\tag{1}
\]

Lemma 1. --- Let X and Y be two locally compact spaces, K (resp. L) a compact subset of X (resp. Y).

\begin{enumerate}
\def\labelenumi{(\roman{enumi})}
\tightlist
\item
  The restriction to \(\mathcal{K}(\mathrm{X}\times \mathrm{Y},\mathrm{K}\times \mathrm{L};\mathbf{C})\) of the canonical bijection
\end{enumerate}

\[
\omega : \mathcal {F} (\mathrm{X} \times \mathrm{Y}; \mathbf {C}) \rightarrow \mathcal {F} \bigl (\mathrm{X}; \mathcal {F} (\mathrm{Y}; \mathbf {C}) \bigr)
\]

(S, R, §4, No.~14) is an isometry of the Banach space \(\mathcal{K}(\mathrm{X}\times\mathrm{Y},\mathrm{K}\times\mathrm{L};\mathbf{C})\) onto the Banach space \(\mathcal{K}(\mathrm{X},\mathrm{K};\mathcal{K}(\mathrm{Y},\mathrm{L};\mathbf{C}))\) .

\begin{enumerate}
\def\labelenumi{(\roman{enumi})}
\setcounter{enumi}{1}
\tightlist
\item
  The vector space \(\mathcal{K}(\mathrm{X},\mathrm{K};\mathbf{C})\otimes_{\mathbf{C}}\mathcal{K}(\mathrm{Y},\mathrm{L};\mathbf{C})\) , identified canonically with a subspace of \(\mathcal{K}(\mathrm{X}\times \mathrm{Y},\mathrm{K}\times \mathrm{L};\mathbf{C})\) (A, II, §7, No.~7, comments following the Cor. of Prop. 15), is dense in this Banach space.
\end{enumerate}

It is immediate that the image under \(\omega\) of \(\mathcal{K}(X\times Y,K\times L;C)\) is contained in \(\mathcal{K}(X,K;\mathcal{K}(Y,L;C))\) . Conversely, if u is a continuous mapping of X into \(\mathcal{K}(Y,L;C)\) , with support contained in K, then the mapping \((x,y)\mapsto\mathbf{u}(x)(y)\) of \(X\times Y\) into C is continuous and has support contained in \(K\times L\) , therefore the restriction of \(\omega\) to \(\mathcal{K}(X\times Y,K\times L;C)\) is a bijection of this space onto \(\mathcal{K}(X,K;\mathcal{K}(Y,L;C))\) ; the fact that this restriction is a Banach space isometry follows from the relation

\[
\sup _ {(x, y) \in \mathrm{K} \times \mathrm{L}} | f (x, y) | = \sup _ {x \in \mathrm{K}} \left(\sup _ {y \in \mathrm{L}} | f (x, y) |\right).
\]

This proves (i); on the other hand the image under \(\omega\) , of

\[
\mathcal {K} (\mathrm{X}, \mathrm{K}; \mathbf {C}) \otimes_ {\mathbf {C}} \mathcal {K} (\mathrm{Y}, \mathrm{L}; \mathbf {C})
\]

identified with a subspace of \(\mathcal{K}(\mathrm{X}\times \mathrm{Y},\mathrm{K}\times \mathrm{L};\mathbf{C})\) , is again the space \(\mathcal{K}(\mathrm{X},\mathrm{K};\mathbf{C})\otimes_{\mathbf{C}}\mathcal{K}(\mathrm{Y},\mathrm{L};\mathbf{C})\) but this time identified canonically with a space of mappings of X into \(\mathcal{K}(\mathrm{Y},\mathrm{L};\mathbf{C})\) (A, II, §7, No.~7, Cor. of Prop. 15); but this subspace of \(\mathcal{K}(\mathrm{X},\mathrm{K};\mathcal{K}(\mathrm{Y},\mathrm{L};\mathbf{C}))\) is known to be dense in the latter space (§1, No.~2, Prop. 5), thus the conclusion of (ii) follows from the fact that the restriction of \(\omega\) is a topological isomorphism.

Having proved the lemma, we now observe that every compact subset of \(X \times Y\) is contained in a product \(K \times L\) , where K (resp. L) is a compact subset of X (resp. Y). It therefore follows from Lemma 1, (ii) that the subspace \(\mathcal{K}(X; \mathbf{C}) \otimes_{\mathbf{C}} \mathcal{K}(Y; \mathbf{C})\) is dense in \(\mathcal{K}(X \times Y; \mathbf{C})\) ; since the relation (1) may also be written as \(\nu(g \otimes h) = \lambda(g)\mu(h)\) for \(g \in \mathcal{K}(X; \mathbf{C})\) , \(h \in \mathcal{K}(Y; \mathbf{C})\) , the uniqueness of \(\nu\) follows at once. To prove the existence of \(\nu\) , we shall make use of the following lemma:

Lemma 2. --- With notations as in Lemma 1, for every function \(f \in \mathcal{K}(\mathrm{X} \times \mathrm{Y}, \mathrm{K} \times \mathrm{L}; \mathbf{C})\) the function

\[
y \mapsto h (y) = \int f (x, y) d \lambda (x)\tag{2}
\]

belongs to \(\mathcal{K}(\mathrm{Y},\mathrm{L};\mathbf{C})\)

For every function \(\mathbf{u} \in \mathcal{K}(\mathrm{X};\mathcal{K}(\mathrm{Y},\mathrm{L};\mathbf{C}))\) , the integral \(\int \mathbf{u}(x)d\lambda (x)\) belongs to \(\mathcal{K}(\mathrm{Y},\mathrm{L};\mathbf{C})\) since the latter is a Banach space ( \(\S 3\) , No.~3, Cor. 1 of Prop. 7); but for \(\mathbf{u} = \omega(f)\) and for every \(y \in \mathrm{Y}\) ,

\[
\left\langle \int \mathbf {u} (x) d \lambda (x), \varepsilon_ {y} \right\rangle = \int \mathbf {u} (x) (y) d \lambda (x) = \int f (x, y) d \lambda (x),
\]

whence the lemma.

Consider, then, for every function \(f \in \mathcal{K}(\mathrm{X} \times \mathrm{Y}; \mathbf{C})\) , the number \(\nu(f) = \mu\left(\int f(x, y) d\lambda(x)\right)\) (which we shall also write \(\int d\mu(y) \int f(x, y) d\lambda(x)\) by an abuse of notation); if K (resp. L) is a compact subset of X (resp. Y), then there exists a number \(a_{K}\) (resp. \(b_{L}\) ) such that, for every function \(g \in \mathcal{K}(\mathrm{X}, \mathrm{K}; \mathbf{C})\) (resp. \(h \in \mathcal{K}(\mathrm{Y}, \mathrm{L}; \mathbf{C})\) ), we have \(|\lambda(g)| \leqslant a_{\mathrm{K}} \|g\|\) (resp. \(|\mu(h)| \leqslant b_{\mathrm{L}} \|h\|\) ). It follows that for every function \(f \in \mathcal{K}(\mathrm{X} \times \mathrm{Y}, \mathrm{K} \times \mathrm{L}; \mathbf{C})\) ,

\[
\left| \int f (x, y) d \lambda (x) \right| \leqslant a _ {K} \| f \|
\]

for every \(y \in Y\) , whence \(|\nu(f)| \leqslant a_{K}b_{L}\|f\|\) . The linear form \(\nu\) on \(\mathcal{K}(X \times Y; \mathbf{C})\) is thus a measure on \(X \times Y\) that obviously satisfies (1), which completes the proof of Th. 1.

DEFINITION 1. --- Given two measures \(\lambda\) , \(\mu\) defined, respectively, on two locally compact spaces X, Y, the product measure of \(\lambda\) by \(\mu\) is defined to be the unique measure \(\nu\) on \(X \times Y\) satisfying the relation (1) for every function \(g \in \mathcal{K}(X; \mathbf{C})\) and every function \(h \in \mathcal{K}(Y; \mathbf{C})\) .

In the proof of Th. 1, the roles of the spaces X and Y may be interchanged; canonically identifying \(Y \times X\) and \(X \times Y\) , we thus define on \(X \times Y\) a measure

\[
f \mapsto \int d \lambda (x) \int f (x, y) d \mu (y)
\]

that again satisfies condition (1). We have thus proved the following theorem:

THEOREM 2. --- Let \(\lambda\) , \(\mu\) be two measures defined, respectively, on two locally compact spaces \(X\) , \(Y\) . For every function \(f\) in \(\mathcal{K}(X \times Y; \mathbf{C})\) , the integral of \(f\) with respect to the product measure \(\nu\) of \(\lambda\) by \(\mu\) has the value

\[
\begin{array}{c} \int f (x, y) d \nu (x, y) = \int d \lambda (x) \int f (x, y) d \mu (y) \\ = \int d \mu (y) \int f (x, y) d \lambda (x). \end{array}\tag{3}
\]

Because of the last formula, the integral of f with respect to the product measure \(\nu\) is usually denoted \(\iint f d\lambda d\mu\) , or \(\iint f d\mu d\lambda\) , or \(\iint f \lambda \mu\) , or \(\iint f \mu \lambda\) , or \(\iint f(x, y) d\lambda(x) d\mu(y)\) , or \(\iint f(x, y) d\mu(y) d\lambda(x)\) , or \(\iint f(x, y) \lambda(x)\mu(y)\) , or \(\iint f(x, y) \mu(y)\lambda(x)\) ; it is said to be the double integral of f with respect to \(\lambda\) and \(\mu\) . With this notation, the formula (3) may be written

\[
\begin{array}{c} \iint f (x, y) d \lambda (x) d \mu (y) = \int d \lambda (x) \int f (x, y) d \mu (y) \\ = \int d \mu (y) \int f (x, y) d \lambda (x). \end{array}\tag{4}
\]

Formula (3) shows that if \(\lambda\) and \(\mu\) are real (resp. positive) measures, then the product measure \(\nu\) is real (resp. positive).

Examples. --- 1) The product measure of the Dirac measures \(\varepsilon_{x}\) ( \(x \in X\) ) and \(\varepsilon_{y}\) ( \(y \in Y\) ) is the Dirac measure \(\varepsilon_{(x,y)}\) .

\begin{enumerate}
\def\labelenumi{\arabic{enumi})}
\setcounter{enumi}{1}
\tightlist
\item
  Let us take \(\mathbf{X} = \mathbf{Y} = \mathbf{R}\) , and for \(\lambda\) and \(\mu\) the Lebesgue measure on \(\mathbf{R}\) (§1, No.~3); their product is called Lebesgue measure on \(\mathbf{R}^2\) ; the integral of a function \(f \in \mathcal{K}(\mathbf{R}^2; \mathbf{C})\) with respect to this measure is denoted \(\iint f(x, y) dx dy\) or \(\iint f(x, y) dy dx\) ; formula (4), for a function that is zero outside a compact rectangle \([a, b] \times [c, d]\) , yields the formula
\end{enumerate}

\[
\int_ {c} ^ {d} d y \int_ {a} ^ {b} f (x, y) d x = \int_ {a} ^ {b} d x \int_ {c} ^ {d} f (x, y) d y
\]

proved in FRV, II, §3, No.~6.

Since Lebesgue measure on \(\mathbf{R}\) is invariant under every translation (§1, No.~3), it follows at once that Lebesgue measure on \(\mathbf{R}^2\) is invariant under every translation of \(\mathbf{R}^2\) .

Remark. --- Let E be a Hausdorff locally convex space, and let f be a mapping in \(\mathcal{K}(X \times Y; E)\) such that \(\mathbf{f}(X \times Y)\) is contained in a complete convex subset C of E. Then, for every \(y \in Y\) , the integral \(\mathbf{h}(y) = \int \mathbf{f}(x, y) d\lambda(x)\) belongs to E (§3, No.~3, Prop. 7); moreover, the function h belongs to \(\widetilde{\mathcal{K}}(Y; E)\) : indeed, for every \(z' \in E'\) we have

\[
\langle \mathbf {h} (y), \mathbf {z} ^ {\prime} \rangle = \int \langle \mathbf {f} (x, y), \mathbf {z} ^ {\prime} \rangle d \lambda (x),
\]

therefore \(y \mapsto \langle \mathbf{h}(y), \mathbf{z}' \rangle\) belongs to \(\mathcal{K}(\mathrm{Y}; \mathbf{C})\) by Lemma 2. The integral \(\int h d\mu\) is therefore defined (and a priori belongs to \(E'^{*}\) ); let us show that

\[
\begin{array}{r l} \iint \mathbf {f} (x, y) d \lambda (x) d \mu (y) & = \int d \mu (y) \int \mathbf {f} (x, y) d \lambda (x) \\ & = \int d \lambda (x) \int \mathbf {f} (x, y) d \mu (y), \end{array}\tag{5}
\]

thus generalizing the formula (4). Indeed, for every \(\mathbf{z}' \in \mathrm{E}'\) we have

\[
\begin{array}{r l} \left\langle \iint \mathbf {f} d \lambda d \mu , \mathbf {z} ^ {\prime} \right\rangle & = \iint \langle \mathbf {f}, \mathbf {z} ^ {\prime} \rangle d \lambda d \mu = \int d \mu \int \langle \mathbf {f}, \mathbf {z} ^ {\prime} \rangle d \lambda \\ & = \int \left\langle \int \mathbf {f} d \lambda , \mathbf {z} ^ {\prime} \right\rangle d \mu = \left\langle \int d \mu \int \mathbf {f} d \lambda , \mathbf {z} ^ {\prime} \right\rangle \end{array}
\]

by (4), whence (5).

